\chapter{Planteamiento del problema}\label{chap:planteamiento_problema}

En este capítulo se concreta el problema abordado en el presente trabajo, definiendo la tarea de predicción prebiopsia, su formulación multietiqueta y el alcance de los análisis realizados. Asimismo, se describen las fuentes de información disponibles, las principales restricciones metodológicas de la cohorte y los recursos empleados durante el desarrollo del proyecto. El objetivo es establecer el marco experimental sobre el que se apoyan los capítulos posteriores.

\section{Descripción general del problema}

El problema abordado en este trabajo consiste en analizar si la información disponible antes de una biopsia pulmonar guiada por tomografía computarizada permite anticipar las complicaciones que pueden producirse tras el procedimiento. Para ello, se estudian distintas representaciones de los pacientes, incluyendo variables clínicas tabulares, volúmenes tridimensionales de TC y características cuantitativas derivadas de las imágenes.

El problema se plantea como una tarea de clasificación multietiqueta, en la que se distinguen las distintas complicaciones y se contempla que un mismo paciente pueda presentar más de una. Esta formulación permite estimar riesgos específicos para cada tipo de complicación, en lugar de limitar el análisis a un único riesgo global.

\subsection{Predicción prebiopsia de complicaciones}

Los modelos estiman el riesgo de complicaciones utilizando exclusivamente la información disponible antes de la biopsia: los datos clínicos preprocedimentales y el estudio de TC torácica de cada paciente. A partir de estas fuentes se analizan, de forma aislada o combinada, variables clínicas, características radiómicas, máscaras de estructuras anatómicas y descriptores geométricos del nódulo y de su relación con el pulmón. Aunque el análisis se realiza retrospectivamente, se excluye cualquier variable obtenida durante o después del procedimiento para reproducir el escenario de aplicación clínica y evitar fugas de información.

\subsection{Formulación multietiqueta}

La aparición de complicaciones no constituye un problema multiclase convencional, ya que un mismo paciente puede presentar más de un evento adverso. En particular, existen casos en los que se registran simultáneamente hemorragia y neumotórax. Por este motivo, el problema se formula como una tarea de clasificación multietiqueta.

Para cada paciente \(i\), la salida se representa mediante el vector binario

\[
\mathbf{y}_i =
\left(
y_i^{H},
y_i^{N}
\right)
\in \{0,1\}^{2},
\]

donde \(y_i^{H}\) indica la presencia de hemorragia y \(y_i^{N}\) la presencia de neumotórax. Ambas complicaciones se representan mediante indicadores binarios separados y no son mutuamente excluyentes, por lo que pueden tomar simultáneamente el valor uno. La ausencia de complicaciones no constituye una etiqueta adicional, sino que corresponde al caso en el que no se ha registrado ninguno de los dos eventos, es decir, \(\mathbf{y}_i=(0,0)\).

El derrame pleural se mantiene en el análisis exploratorio para describir íntegramente la información clínica disponible, pero no se incorpora como objetivo de los modelos debido a que solo está presente en un paciente.

Esta formulación permite obtener una predicción diferenciada para cada complicación. No obstante, también introduce dificultades adicionales respecto a una única tarea binaria, como el desbalance entre etiquetas, la escasez de algunas combinaciones y la necesidad de utilizar métricas que evalúen tanto el rendimiento global como el comportamiento específico para cada complicación.

\subsection{Alcance del estudio: predicción y caracterización}

El alcance del trabajo comprende dos perspectivas complementarias. La primera es predictiva y busca evaluar hasta qué punto las distintas fuentes de información disponibles permiten anticipar la aparición de hemorragia, neumotórax o ausencia de complicaciones. Para ello, se comparan modelos basados en variables clínicas, características radiómicas, volúmenes tridimensionales de TC y combinaciones de varias fuentes de información.

La segunda perspectiva es descriptiva y está orientada a comprender mejor la estructura del problema. En este sentido, el estudio no pretende limitarse a obtener una puntuación de rendimiento, sino también analizar qué variables o representaciones aportan información relevante, cómo afecta la posible incertidumbre de las etiquetas y qué perfiles de pacientes presentan una frecuencia diferencial de complicaciones. Para abordar estos aspectos se incorporan técnicas como el \textit{PU learning}, destinado a explorar escenarios con posibles positivos no identificados, y el descubrimiento de subgrupos, utilizado para obtener reglas interpretables que caractericen grupos concretos de pacientes.

Por tanto, la finalidad del trabajo no consiste exclusivamente en desarrollar el modelo con mayor rendimiento, sino en realizar una comparación metodológica amplia que permita valorar la utilidad de las distintas fuentes de información y comprender las limitaciones asociadas a una cohorte clínica de tamaño reducido y heterogénea.



\section{Descripción de los datos disponibles}
\label{sec:descripcion_datos}

El conjunto de datos utilizado en este trabajo fue proporcionado de forma anonimizada por el Hospital Universitario Clínico San Cecilio de Granada. La información disponible combina datos clínicos tabulares y volúmenes de tomografía computarizada.

La cohorte inicial está formada por 210 pacientes con información clínica. La disponibilidad y calidad de los estudios de imagen y de sus representaciones derivadas no son necesariamente uniformes para todos los casos.

\subsection{Datos clínicos tabulares}\label{sec:variables_disponibles}

La información clínica se proporcionó mediante un fichero tabular en el que cada fila corresponde a un paciente. El conjunto original contiene variables demográficas, antecedentes clínicos, información anatomopatológica y datos relacionados con la aparición de complicaciones. Entre las variables disponibles se encuentran:

\begin{itemize}[itemsep=2pt, topsep=4pt]
    \item \textbf{Identificador del paciente}: código empleado para relacionar la información clínica con los estudios de imagen y sus archivos derivados.

    \item \textbf{Sexo y edad}: variables demográficas del paciente.

    \item \textbf{Tipo de cáncer}: subtipo tumoral identificado tras el estudio anatomopatológico.

    \item \textbf{Complicación}: indicador global de la aparición o ausencia de alguna complicación asociada a la biopsia.

    \item \textbf{Tipo de complicación}: especificación de los eventos registrados, principalmente hemorragia, neumotórax y, de forma excepcional, derrame pleural.

    \item \textbf{Factores de riesgo}: antecedentes y condiciones clínicas potencialmente relevantes, como tabaquismo, hipertensión arterial, diabetes mellitus u otras comorbilidades.

    \item \textbf{Patologías pulmonares}: enfermedades respiratorias previas o alteraciones pulmonares, como enfisema, fibrosis, bronquiectasias o hipertensión pulmonar.

    \item \textbf{TC}: dispositivo utilizado para la adquisición del estudio de tomografía computarizada.

    \item \textbf{Extensión a los 1-5 años}: información de seguimiento relativa a la evolución o extensión de la enfermedad durante dicho periodo.

    \item \textbf{Anticoagulación}: información sobre el tratamiento anticoagulante del paciente.

    \item \textbf{Mutación}: presencia de alteraciones moleculares o marcadores genéticos registrados durante el proceso diagnóstico.
\end{itemize}

% \subsection{Estudios de TC en formato DICOM}

% Los estudios de imagen originales se proporcionaron en formato DICOM (\textit{Digital Imaging and Communications in Medicine}), estándar utilizado habitualmente para el almacenamiento y el intercambio de imágenes médicas. Cada archivo DICOM contiene la matriz de intensidades correspondiente a una imagen y un encabezado con metadatos técnicos y espaciales, como la modalidad, el fabricante, la resolución, el espaciado entre píxeles, la posición del corte o el identificador de la serie.

% Tras revisar los datos disponibles, se seleccionó para cada paciente el estudio de TC principal adecuado para el análisis, obteniéndose 209 estudios válidos. Los volúmenes originales presentan heterogeneidad en el número de cortes, las dimensiones en el plano, el espaciado físico y los parámetros de adquisición. Estas propiedades se analizan de forma detallada en el Capítulo~\ref{chap:eda}, donde también se justifican las operaciones posteriores de homogeneización y control de calidad.

% \subsection{Volúmenes NIfTI derivados}

% A partir de las series DICOM seleccionadas se generaron volúmenes en formato NIfTI (\textit{Neuroimaging Informatics Technology Initiative}). Este formato permite almacenar en un único archivo la matriz tridimensional de intensidades y la información espacial necesaria para localizar físicamente sus vóxeles, por lo que facilita el uso de herramientas de procesamiento, segmentación y aprendizaje automático.

% Se dispuso de 206 volúmenes NIfTI válidos. La diferencia respecto a los 209 estudios DICOM seleccionados se debe a la existencia de algunos casos con información insuficiente o problemas de geometría que impedían obtener un volumen tridimensional adecuado. Estos casos se identificaron durante el análisis de calidad y no se utilizaron en los experimentos que requerían representación volumétrica.

% Los archivos NIfTI constituyen la representación de referencia para la aplicación de segmentaciones anatómicas, la extracción de características radiómicas y el entrenamiento de modelos basados en imagen. Las transformaciones específicas aplicadas sobre estos volúmenes, como normalización de intensidades, ajuste espacial o redimensionado, se describen posteriormente en la sección de preprocesamiento.

\subsection{Datos de imagen}

Los estudios de TC se han proporcionado originalmente en formato DICOM (\textit{Digital Imaging and Communications in Medicine}), estándar utilizado para almacenar imágenes médicas junto con sus metadatos técnicos y espaciales. Estos metadatos incluyen información sobre la adquisición y la geometría de las imágenes, necesaria para reconstruir los estudios como volúmenes tridimensionales.

Para facilitar su procesamiento, los estudios DICOM se han convertido a volúmenes en formato NIfTI (\textit{Neuroimaging Informatics Technology Initiative}). Este formato integra en un único archivo la matriz tridimensional de intensidades y su información espacial, lo que facilita el procesamiento computacional de las imágenes y su correspondencia con otras representaciones volumétricas.

\subsection{Máscaras anatómicas y segmentaciones}

Además de los volúmenes de TC, se dispuso de máscaras tridimensionales de distintas estructuras anatómicas relevantes para el estudio. Estas máscaras delimitan los vóxeles pertenecientes a una región concreta y permiten restringir el análisis a zonas potencialmente relacionadas con el procedimiento de biopsia.

Las principales estructuras segmentadas fueron:

\begin{itemize}[itemsep=2pt, topsep=4pt]
\item \textbf{Pulmón}: delimita el parénquima pulmonar y permite excluir gran parte de las estructuras externas al área principal de estudio.
\item \textbf{Nódulo}: identifica la lesión pulmonar sobre la que se realiza la biopsia y constituye una de las principales regiones de interés del análisis radiómico.
\item \textbf{Vasos pulmonares}: representa las estructuras vasculares presentes en el pulmón y permite estudiar información anatómica potencialmente relacionada con el riesgo hemorrágico.
\item \textbf{Estructuras bronquiales}: recoge la vía aérea visible y proporciona información adicional sobre el entorno anatómico de la lesión.
\end{itemize}

Las máscaras se almacenaron en formato NIfTI y conservaron la correspondencia espacial con el volumen de TC de cada paciente. Se utilizaron para extraer características radiómicas y medidas geométricas, así como para construir distintas representaciones de entrada para los modelos tridimensionales. Su preparación y su utilización experimental se describen en los capítulos correspondientes.

% \subsection{Correspondencia entre pacientes, imágenes y etiquetas}

% La integración de las distintas fuentes se realizó mediante el identificador anónimo asignado a cada paciente. Este código permitió relacionar la fila correspondiente del fichero clínico con la serie DICOM seleccionada, el volumen NIfTI, las máscaras anatómicas y las etiquetas de complicación.

% La cobertura clínica fue completa para los 210 pacientes, mientras que 209 disponían de una TC principal válida y 206 contaban con un volumen NIfTI adecuado para el procesamiento tridimensional. En consecuencia, se definieron diferentes cohortes experimentales según la información requerida en cada análisis. Los experimentos exclusivamente clínicos pudieron utilizar la cohorte completa, mientras que los modelos radiómicos, geométricos o basados directamente en imagen se limitaron a los pacientes con los archivos necesarios.

% Esta separación se tuvo en cuenta al comparar configuraciones. Cuando el objetivo era estudiar la aportación relativa de dos fuentes de información, se utilizaron cohortes emparejadas formadas por los mismos pacientes, evitando atribuir a una modalidad diferencias que pudieran deberse únicamente a cambios en la composición de la muestra.


% \section{Legalidad de los datos}

% Los datos médicos empleados en este proyecto corresponden a volúmenes de tomografía computarizada en formato DICOM, previamente anonimizados por el equipo clínico responsable. Este proceso consistió en la eliminación de los identificadores personales directos e indirectos contenidos en los metadatos de los archivos DICOM, de acuerdo con las directrices recogidas en el \emph{Considerando 26 del Reglamento (UE) 2016/679 (RGPD)} \cite{gdprRecital26}, según el cual los datos anonimizados de forma irreversible quedan excluidos del ámbito de aplicación de dicho reglamento.

% Asimismo, conforme a lo establecido en la \emph{Ley Orgánica 3/2018}, de \emph{Protección de Datos Personales y garantía de los derechos digitales}, en particular en su \emph{Disposición adicional decimoséptima, apartado 2, letra e} \cite{lopd2018}, se considera que los datos utilizados no presentan una probabilidad razonable de reidentificación de los pacientes. De este modo, el tratamiento de la información empleada en este estudio se ajusta a la normativa vigente en materia de protección de datos personales y sigue las recomendaciones establecidas por la \emph{Agencia Española de Protección de Datos (AEPD)} \cite{aepdAnonimizacion}.

% La correcta anonimización de los datos constituye un requisito fundamental para el desarrollo responsable de proyectos de inteligencia artificial en el ámbito sanitario, ya que permite preservar la privacidad de los pacientes y garantizar el respeto a sus derechos fundamentales durante todo el proceso de análisis y modelado.
\section{Legalidad de los datos}

Los volúmenes de TC utilizados fueron previamente anonimizados por el equipo clínico responsable mediante la eliminación de los identificadores personales contenidos en los metadatos DICOM. El tratamiento de estos datos se ha planteado de acuerdo con el \emph{Considerando 26 del Reglamento (UE) 2016/679, General de Protección de Datos (RGPD)} \cite{gdprRecital26}, la \emph{Disposición adicional decimoséptima de la Ley Orgánica 3/2018, de Protección de Datos Personales y garantía de los derechos digitales} \cite{lopd2018}, y las recomendaciones de la \emph{Agencia Española de Protección de Datos (AEPD)} \cite{aepdAnonimizacion}, con el objetivo de preservar la privacidad de los pacientes y minimizar el riesgo de reidentificación.



% \section{Restricciones y dificultades metodológicas}

% El problema estudiado presenta varias limitaciones derivadas de la naturaleza clínica de los datos, su dimensionalidad y la forma en que se definen y observan las etiquetas de complicación. Estas restricciones condicionan tanto el diseño experimental como la interpretación de los resultados y deben tenerse en cuenta al valorar la capacidad de generalización de los modelos.

% \subsection{Tamaño muestral y distribución de las etiquetas}

% La principal limitación del estudio es el tamaño reducido de la cohorte en comparación con la complejidad de los modelos evaluados. Esta dificultad es especialmente relevante en los enfoques de aprendizaje profundo, donde el número de parámetros puede ser elevado respecto al número de pacientes disponibles, aumentando el riesgo de sobreajuste.

% La formulación multietiqueta introduce además un desequilibrio entre los distintos resultados clínicos. Aunque la etiqueta global binaria de complicación presenta una distribución equilibrada, las complicaciones específicas tienen un soporte desigual y algunas combinaciones de etiquetas aparecen en muy pocos pacientes. Por este motivo, el derrame pleural no se incluyó como objetivo predictivo independiente, al estar registrado únicamente en un paciente.

% También debe considerarse la posible incertidumbre de las etiquetas clínicas. La ausencia de una complicación registrada no garantiza necesariamente que el paciente constituya un negativo completamente fiable, ya que pueden existir eventos leves, registros incompletos o diferencias en los criterios de anotación. Esta posibilidad motivó la exploración de técnicas de \textit{Positive-Unlabeled Learning}.

% \subsection{Heterogeneidad y alta dimensionalidad de los datos}

% Los estudios de TC proceden de diferentes equipos y protocolos de adquisición, por lo que presentan variabilidad en el número de cortes, la resolución, el espaciado entre vóxeles, el grosor de corte y los parámetros de reconstrucción. Esta heterogeneidad puede introducir diferencias técnicas no relacionadas con la aparición de complicaciones y dificultar el aprendizaje de patrones generalizables.

% Asimismo, tanto los volúmenes de TC como las representaciones radiómicas presentan una dimensionalidad elevada. En el caso de la radiómica, el número de características puede ser comparable o superior al número de pacientes, lo que incrementa el riesgo de redundancia y sobreajuste. Por ello, fue necesario aplicar procedimientos de homogeneización, reducción de dimensionalidad y selección de características dentro del proceso de validación.

% \subsection{Riesgo de fuga de información}

% Una dificultad metodológica central es evitar que los modelos utilicen información que no estaría disponible antes de la biopsia. Por este motivo, se excluyeron como predictores las variables obtenidas después del procedimiento, durante el diagnóstico o en el seguimiento posterior.

% Además, todas las transformaciones dependientes de los datos, como la imputación, el escalado, la selección de características o la reducción de dimensionalidad, deben ajustarse exclusivamente con los pacientes del conjunto de entrenamiento de cada partición. Aplicarlas previamente sobre la cohorte completa produciría una fuga de información y una estimación optimista del rendimiento.

% El identificador anónimo del paciente tampoco se utilizó como predictor, ya que su estructura contenía información codificada sobre determinadas características clínicas y sobre el resultado de la biopsia.

% % \subsection{Validación y generalización clínica}

% % La disponibilidad de datos procedentes de un único centro impide realizar una validación externa independiente. En consecuencia, la evaluación se basa principalmente en validación cruzada, procurando mantener particiones coherentes entre experimentos y evitar que datos del mismo paciente aparezcan simultáneamente en entrenamiento y evaluación.

% % Los resultados obtenidos deben interpretarse como evidencia preliminar sobre la viabilidad de las estrategias estudiadas. Antes de considerar una aplicación clínica sería necesario validar los modelos sobre cohortes externas, procedentes de otros centros y equipos de adquisición, y realizar una evaluación prospectiva con participación del equipo médico. Por tanto, el trabajo no pretende proporcionar una herramienta clínica definitiva, sino identificar enfoques prometedores, limitaciones metodológicas y posibles líneas de mejora.

\section{Restricciones y dificultades metodológicas}

El diseño experimental y la interpretación de los resultados están condicionados por las siguientes características del problema:

\begin{itemize}
    \item \textbf{Tamaño muestral reducido:} el número de pacientes es limitado en relación con la complejidad de algunos modelos, especialmente los de \textit{deep learning}, lo que incrementa el riesgo de sobreajuste y limita su capacidad de generalización.

    \item \textbf{Distribución e incertidumbre de las etiquetas:} la hemorragia y el neumotórax presentan soportes desiguales y algunas combinaciones son poco frecuentes. El derrame pleural se consideró únicamente en el análisis exploratorio. Además, la posible existencia de eventos no registrados motivó la exploración de técnicas de \textit{PU learning}.

    \item \textbf{Heterogeneidad de las imágenes:} los estudios de TC proceden de distintos equipos y protocolos, con variaciones en la resolución, el espaciado, el grosor de corte y los parámetros de reconstrucción. Estas diferencias pueden introducir variabilidad técnica no relacionada con las complicaciones.

    \item \textbf{Alta dimensionalidad:} los volúmenes de TC y las características radiómicas contienen un número elevado de variables respecto al tamaño de la cohorte. Para reducir la redundancia y el riesgo de sobreajuste se aplicaron procedimientos de homogeneización, reducción de dimensionalidad y selección de características dentro del proceso de validación.

    \item \textbf{Literatura específica limitada:} como se expone en el Capítulo~\ref{chap:estado_del_arte}, la evidencia sobre la predicción de complicaciones de la biopsia pulmonar mediante inteligencia artificial continúa siendo limitada y heterogénea. La mayoría de los trabajos se centra en una única complicación, especialmente el neumotórax, o emplea variables clínicas y procedimentales tabulares. Son menos frecuentes los estudios que incorporan directamente la imagen de TC, la radiómica o una formulación multietiqueta. Esta escasez de estudios comparables dificulta la selección de referencias experimentales y la existencia de protocolos de evaluación consolidados.

\end{itemize}


\section{Metodología}

La metodología de trabajo se enmarca en un enfoque empírico y aplicado, situado en la intersección entre la ingeniería informática, la inteligencia artificial y el ámbito médico.

El desarrollo del trabajo se ha organizado siguiendo un procedimiento estructurado, inspirado en el método científico y adaptado a las particularidades del problema clínico estudiado. Las fases principales de la metodología seguida son las siguientes:

\begin{itemize}
    \item \textbf{Observación.} Se estudió el contexto clínico de las biopsias pulmonares, las complicaciones asociadas y el estado del arte sobre inteligencia artificial aplicada a imagen médica y predicción de riesgos clínicos.

    \item \textbf{Formulación de hipótesis.} Se planteó que la aparición de complicaciones podría estar relacionada con patrones presentes en los datos clínicos tabulares y en los volúmenes de TC, por lo que se evaluaron modelos capaces de explotar ambas fuentes de información.

    \item \textbf{Recogida de datos.} Se trabajó con datos anonimizados proporcionados por el Hospital Universitario Clínico San Cecilio de Granada, incluyendo estudios de TC en formato DICOM y variables clínicas asociadas a cada paciente.

    \item \textbf{Preprocesamiento y segmentación.} Se aplicaron procedimientos específicos para preparar los datos: limpieza y codificación de variables clínicas, selección y conversión de series de TC, normalización de intensidades, homogeneización volumétrica y segmentación anatómica.

    \item \textbf{Prueba de hipótesis.} Se entrenaron y evaluaron diferentes aproximaciones basadas en aprendizaje automático, \textit{deep learning}, estrategias multimodales y técnicas más avanzadas de IA para analizar su capacidad predictiva.

    \item \textbf{Comprobación y análisis.} Los resultados se estudiaron mediante métricas de clasificación e interpretación técnica y clínica, considerando las limitaciones del conjunto de datos y la influencia del preprocesamiento.

    \item \textbf{Conclusiones.} Finalmente, se extrajeron conclusiones sobre la viabilidad del enfoque propuesto, la utilidad de las fuentes de información disponibles y las posibles líneas futuras de mejora y validación clínica.
\end{itemize}





\section{\textit{Software}}

El desarrollo se ha realizado principalmente en Python 3.12.11. Para la implementación y el entrenamiento de los modelos de \textit{deep learning} se han utilizado PyTorch 2.5.1 y MONAI 1.4.0 \cite{cardoso2022monai}. El procesamiento de las imágenes médicas se ha llevado a cabo mediante NiBabel 5.3.2, Pydicom 3.0.1 y SimpleITK 2.4.1, mientras que las segmentaciones anatómicas se han generado con TotalSegmentator \cite{wasserthal2023totalsegmentator}. La extracción de características radiómicas se ha realizado mediante PyRadiomics \cite{VanGriethuysen2017}.

Para el tratamiento de los datos y la implementación de los modelos de aprendizaje automático se han empleado NumPy 2.2.6, Pandas 2.2.3, scikit-learn 1.5.2 y XGBoost 1.7.6. Los análisis exploratorios y las visualizaciones se han desarrollado con Jupyter Notebooks, Matplotlib y Seaborn. Asimismo, se han utilizado Visual Studio Code y Git para el desarrollo y el control de versiones.

El código desarrollado se encuentra disponible en el repositorio de GitHub \url{https://github.com/mcribi/Prediction-Complication-Lung-Biopsies-AI}. Los datos clínicos y las imágenes médicas no se han incluido en el repositorio debido a su carácter sensible.


\section{\textit{Hardware}}

Para el desarrollo experimental del proyecto se utilizaron los servidores de cómputo de alto rendimiento NGPU proporcionados por la Universidad de Granada (UGR), ubicados en el CPD Santa Lucía. Esta infraestructura permitió ejecutar las tareas con mayor coste computacional, especialmente el entrenamiento y la evaluación de los modelos.

La infraestructura NGPU de la UGR dispone de un clúster de nodos de cómputo con distintas configuraciones de GPU, entre las que se incluyen GTX Titan X Pascal, GTX Titan Xp, RTX 2080 Ti, Titan RTX, Tesla V100 de 32 GB y Tesla A100 de 40 GB. Los experimentos se ejecutaron en esta infraestructura mediante Slurm, utilizado para gestionar las colas de trabajo y la asignación de recursos entre los distintos nodos de cómputo.


\section{Presupuesto}

A continuación, se presenta el presupuesto estimado para la realización del proyecto. Este presupuesto incluye los recursos humanos necesarios para el desarrollo del trabajo, así como una estimación del coste asociado al uso de infraestructura de cómputo para el entrenamiento y evaluación de los modelos. La duración total del proyecto se estima en \textit{6 meses}, desde marzo de 2026 hasta agosto de 2026, ambos incluidos.

Se excluyen de este presupuesto los costes personales propios del investigador, como el uso de ordenador personal, alquiler de vivienda o consumo eléctrico doméstico. Asimismo, dado que las tareas de cómputo se han realizado en servidores proporcionados por la Universidad de Granada, los costes de infraestructura se consideran como una estimación equivalente y no como un gasto directo del proyecto.

Las fuentes utilizadas para la estimación de costes son \textit{Glassdoor}, para la estimación del coste asociado al perfil de \textit{Junior Software Engineer}\footnote{\url{https://www.glassdoor.es/Sueldos/junior-software-engineer-sueldo-SRCH_KO0,24.htm}} y al perfil médico\footnote{\url{https://www.glassdoor.es/Sueldos/medico-sueldo-SRCH_KO0,6.htm}}, y \textit{Google Cloud}, para la estimación equivalente del coste de infraestructura de cómputo con GPU\footnote{\url{https://cloud.google.com/compute/gpus-pricing}}.
\subsection{Recursos humanos}

El equipo del proyecto está compuesto por el investigador principal, responsable del desarrollo técnico, análisis de datos y evaluación de modelos; un médico colaborador, encargado de la orientación clínica y supervisión del problema médico; y el tutor académico, responsable de la dirección y seguimiento del trabajo. En la Tabla~\ref{tab:recursos-humanos} se detallan los costes asociados.

\begin{table}[!htbp]
    \centering
    \begin{tabularx}{\textwidth}{|X|c|c|c|}
        \hline
        \textbf{Recurso} & \textbf{Coste mensual} & \textbf{Duración} & \textbf{Coste total} \\ \hline
        Investigador principal & 1.800 € & 6 meses & 10.800 € \\
        Médico colaborador & 3.000 € & 6 meses & 18.000 € \\
        Tutor académico (25 h) & 2.000 € & 1 mes & 2.000 € \\ \hline
        \textbf{Total recursos humanos} & & & \textbf{30.800 €} \\ \hline
    \end{tabularx}
    \caption{Presupuesto estimado destinado a recursos humanos.}
    \label{tab:recursos-humanos}
\end{table}

\subsection{Capacidad de cómputo}

El entrenamiento y evaluación de modelos sobre volúmenes de tomografía computarizada requiere infraestructura de alto rendimiento, especialmente en los experimentos con modelos tridimensionales. En este proyecto se han utilizado los servidores NGPU proporcionados por la Universidad de Granada, por lo que no se ha generado un coste directo asociado a servicios comerciales en la nube.

No obstante, para reflejar el valor aproximado de la infraestructura empleada, se incluye una estimación equivalente basada en el uso de servidores con GPU y almacenamiento durante el periodo de desarrollo experimental. Esta estimación se muestra en la Tabla~\ref{tab:computo}.

\begin{table}[!htbp]
    \centering
    \begin{tabularx}{\textwidth}{|X|c|c|c|}
        \hline
        \textbf{Recurso} & \textbf{Coste mensual} & \textbf{Duración} & \textbf{Coste total} \\ \hline
        Servidores GPU & 100 € & 6 meses & 600 € \\
        Almacenamiento (1 TB) & 50 € & 6 meses & 300 € \\ \hline
        \textbf{Total capacidad de cómputo} & & & \textbf{900 €} \\ \hline
    \end{tabularx}
    \caption{Estimación equivalente del coste asociado a capacidad de cómputo.}
    \label{tab:computo}
\end{table}

\subsection{Materiales y otros gastos}

No se prevén gastos adicionales en licencias de \textit{software}, ya que las herramientas utilizadas durante el desarrollo del proyecto son de código abierto o de uso libre en el ámbito académico, como Python, PyTorch, MONAI, PyRadiomics, scikit-learn y otras bibliotecas empleadas para el análisis y modelado de los datos. Tampoco se consideran necesarias compras de material adicional.

\subsection{Presupuesto total estimado}

El presupuesto total estimado para el desarrollo del proyecto durante un periodo de 6 meses asciende a \textbf{31.700 €}. Este importe incluye los recursos humanos y una estimación equivalente del uso de infraestructura de cómputo. Si se considera únicamente el gasto directo, excluyendo el coste imputado de los servidores NGPU cedidos por la Universidad de Granada, el presupuesto sería de \textbf{30.800 €}.
